% Separate analytical increment. Do not merge without explicit approval.
% Label prefix: inc:LPR:v1:
\documentclass[11pt]{article}
\usepackage[T1]{fontenc}
\usepackage[utf8]{inputenc}
\usepackage{lmodern,amsmath,amssymb,amsthm,mathtools,microtype}
\usepackage[margin=0.95in]{geometry}
\usepackage{booktabs,array,xurl}
\usepackage[hidelinks]{hyperref}
\usepackage{fancyhdr}
\pagestyle{fancy}\fancyhf{}
\fancyhead[L]{\small Learned limiting reversal and the retention interface}
\fancyhead[R]{\small Review increment v1}
\fancyfoot[C]{\thepage}
\setlength{\headheight}{14pt}
\setlength{\emergencystretch}{3em}
\allowdisplaybreaks[2]
\numberwithin{equation}{section}
\newtheorem{theorem}{Theorem}[section]
\newtheorem{proposition}[theorem]{Proposition}
\newtheorem{lemma}[theorem]{Lemma}
\newtheorem{corollary}[theorem]{Corollary}
\theoremstyle{remark}\newtheorem{remark}[theorem]{Remark}
\newcommand{\E}{\mathbb E}
\newcommand{\R}{\mathbb R}
\newcommand{\dd}{\,d}
\newcommand{\norm}[1]{\left\lVert#1\right\rVert}
\DeclareMathOperator{\diag}{diag}
\title{Learned reversal, limiting persistence, and an analytical retention interface}
\author{Separate proof increment for independent review}
\date{Version 1 --- October 3, 2026}
\begin{document}\maketitle

\begin{abstract}
We combine the common learned witnesses of AN06 v2 with the selected-orbit
permanence theorem on the same parameter--probe rectangle. We also prove,
for the full five-dimensional selected orbit rather than its fully learned
truncation, that the weak scaled coordinates grow only as
$\sqrt{\log(2+s)}$. The proof retains the exponentially decaying weak-mass
transient. Finally, we state an exact fixed-label mass-retention inequality
for the original population and a conditional seed--shape exponent
calculation. These last results do not prove capture or isotropic entry.
The initialized Theorems A and B, canonical finite-noise coverage, and an
explicit growing-horizon transfer remain open. No numerical trajectory,
interval arithmetic, or numerical sign test is used as a premise.
\end{abstract}

\section{Contracts and the selected orbit}\label{inc:LPR:v1:contract}
This is a new, separate file. The checkpoint, foundations, AN02, AN06, and
linear/persistence increment are unchanged. The results below are not an
integration approval and do not replace the initialized proof obligations.

The imported analytical results are:
\begin{itemize}
\item \cite{P}, \texttt{pa:branch} and \texttt{pa:limitreversal}: the selected
coherent orbit, global forward existence, its source identities,
$\max(|x|,|y|,|u|,|v|)\le C(1+s)$ for $s\ge0$, and the required invariants;
\item \cite{LP}: the permanence, slow-growth, and uniform-clearing
results (labels \path{inc:LP:v1:permanence},
\path{inc:LP:v1:monotonic-theorem}, and
\path{inc:LP:v1:uniform-clear}). They supply eventual gate inactivity,
$x=O(\sqrt{\log(2+s)})$, bounded $v-x$, and uniform clearing on compact
parameter sets;
\item \cite{W}, \texttt{inc:AN06:v2:uniformtheorem}: learned witnesses on
its specified closed rectangle $\mathcal P_*$ around $(2/3,5\pi/9)$, with
positive constants defined analytically from the selected orbit.
\end{itemize}
The provenance of these results is their displayed analytical proofs, not
their numerical evaluation. They remain subject to independent review.

Let $0<\rho<1$, $\lambda=\rho^2$, and let $\phi,\Phi$ denote the standard
normal density and distribution function. Put
\[
 K(z)=\phi(z)+z\Phi(z),\qquad g(z)=K(z)-z>0.
\]
The normalized time is $\tau=\mu_1^2t_{\rm phys}$. The leading weak
half-mass-centered clock is $s$, so
$n(s)=(1+e^{-\lambda s})^{-1}$. The selected strong-learned coherent orbit
obeys
\begin{align}
 x'&=\frac\lambda2[nv-x+(1-n)y]\Phi(\rho x),\nonumber\\
 y'&=\frac12[-y-nK(u)+\rho(1-n)K(\rho x)],\nonumber\\
 u'&=\frac12[(1-n)(\phi(u)-\lambda u)-(nu+y)\Phi(u)],\nonumber\\
 v'&=\frac12[\rho K(\rho x)-\lambda v].
 \label{inc:LPR:v1:orbit}
\end{align}
It is the analytically specified branch incoming from the resident saddle,
not a trajectory fitted to empirical coordinates. The strong leading mass
is identically one. For the probe offset $\zeta$, its observable is
\begin{equation}\label{inc:LPR:v1:error}
 \mathcal E_\zeta(s)=\frac12(\zeta-x(s))_+^2
 -(\zeta-x(s))_+(\zeta-y(s)).
\end{equation}
These are statements within the defined limiting system. They do not yet
assert a uniform approximation to the original positive-noise population.

\section{One limiting theorem on the AN06 rectangle}
\label{inc:LPR:v1:combined}
Retain the exact notation of AN06 v2. Its rectangle
$\mathcal P_*=[\rho_c-\delta_\rho,\rho_c+\delta_\rho]
\times[\zeta_c-\delta_\zeta,\zeta_c+\delta_\zeta]$ has positive side
lengths, center $(\rho_c,\zeta_c)=(2/3,5\pi/9)$, and is contained in
$[13/20,7/10]\times[87/50,7/4]$. Let $p_*,\gamma_0,r,\sigma>0$ and
\[
 J=[s_*-r,s_*+r],\quad
 J_-=[s_*-r,s_*-r/2],\quad
 J_+=[s_*+r/2,s_*+r]
\]
be its common analytical witnesses. The symbol $\sigma$ here is the AN06
mass slack, not a data-noise parameter; normalized data noise is $q$.

\begin{theorem}[Learned signed reversal followed by eventual zero benefit]
\label{inc:LPR:v1:combinedtheorem}
On $\mathcal P_*$, the selected coherent leading orbit has strong mass one,
weak mass $n\ge1/4+\sigma/2$ on $J$, and
\begin{align}
 d_1&\le-5p_*/8, &d_2&\ge5p_*/8 &&\text{on }J,\nonumber\\
 \mathcal E_\zeta'&\le-3\gamma_0/4&&&\text{on }J_-,\nonumber\\
 \mathcal E_\zeta'&\ge3\gamma_0/4&&&\text{on }J_+.
 \label{inc:LPR:v1:combinedmargins}
\end{align}
There is an interior minimizer $s_{\min}$ of $\mathcal E_\zeta$ on $J$,
possibly parameter dependent, with
\[
 \mathcal E_\zeta(s_*-r)-\mathcal E_\zeta(s_{\min})
 \ge\frac{3\gamma_0r}{8},\qquad
 \mathcal E_\zeta(s_*+r)-\mathcal E_\zeta(s_{\min})
 \ge\frac{3\gamma_0r}{8}.
\]
There is also one finite $S>s_*+r$, independent of $(\rho,\zeta)\in
\mathcal P_*$, such that
\[
 \mathcal E_\zeta(s)=0\qquad(s\ge S).
\]
The parameter widths, margins, and clearing time are analytical constants,
not decimal evaluations of a saved orbit.
\end{theorem}
\begin{proof}
The mass and source-rate assertions are precisely AN06 v2,
\texttt{uniformtheorem}. Each signed subinterval has length $r/2$.
Integrating gives the two stated drops relative to the minimum; the same
inequalities exclude the endpoints of $J$ as minimizers. The selected
branch is continuous on $J$, so a minimizer exists.
Apply \cite{LP}, \texttt{uniform-clear}, to the compact projection of
$\mathcal P_*$ on the $\rho$ axis, using its largest offset as $\zeta_+$.
AN06's branch-parameter continuity and the uniform forward growth estimate
supply that proposition's hypotheses. Enlarge the resulting time beyond
$s_*+r$. After it, $x_\rho>\zeta_+$ for every parameter, and both positive
parts in \eqref{inc:LPR:v1:error} vanish.
\end{proof}

\begin{remark}[The range is not enlarged by combining results]
Permanence alone holds for every fixed $0<\rho<1$ in the selected model.
The combined learned-witness statement uses the smaller AN06 rectangle.
A quarter-mass statement in this leading system is not itself a
finite-noise response/loss theorem: AN06's separate static learning lemma
has additional hypotheses. This theorem does not cover a non-small
bridge, prove initialized branch selection, or certify the canonical
$(q,S_0)=(0.05,2\cdot10^{-4})$ trajectory.
\end{remark}

\section{Weak-coordinate growth in the full selected system}
\label{inc:LPR:v1:weakgrowth}
The full selected system, not the substitution $n=1$, admits a useful
late-time comparison. Set
\[
 c=-y,\qquad w=c-u,\qquad Q=\Phi(u),\qquad r_n=1-n.
\]
We reserve $r_n$ for the mass deficit to distinguish it from the AN06
window radius. Rearranging \eqref{inc:LPR:v1:orbit} gives exactly
\begin{align}
 u'&=\tfrac12wQ+f, &
 f&=\tfrac{r_n}{2}[K(u)-\lambda u],\nonumber\\
 c'&=\tfrac12[g(u)-w]-e, &
 e&=\tfrac{r_n}{2}[K(u)+\rho K(\rho x)],\nonumber\\
 w'&=\tfrac12g(u)-a(s)w-h, &
 a(s)&=\tfrac12(1+Q),\quad h=e+f.
 \label{inc:LPR:v1:weakpair}
\end{align}
For $0<\lambda<1$, $f,e,h$ are positive at finite states. Indeed
$K(u)-\lambda u>0$ both for $u\ge0$ and $u<0$.
The imported linear growth bound and $r_n\le e^{-\lambda s}$ imply
\begin{equation}\label{inc:LPR:v1:weakforcing}
 0\le h(s)\le C_h(1+s)e^{-\lambda s}\quad(s\ge0),
 \qquad \tfrac12\le a(s)\le1.
\end{equation}

\begin{proposition}[Weak divergence, logarithmic growth, and eventual drift]
\label{inc:LPR:v1:weakgrowththeorem}
On the selected orbit, for every fixed $0<\rho<1$,
\[
 u(s)\longrightarrow+\infty,\qquad c(s)=-y(s)\longrightarrow+\infty,
\]
\[
 |c(s)-u(s)|\le C_w\quad(s\ge0),\qquad
 |u(s)|+|y(s)|=O(\sqrt{\log(2+s)}).
\]
Moreover, $c(s)-u(s)>0$ and $u'(s)>0$ for all sufficiently large $s$.
The constants and starting times can depend on $\rho$. These conclusions
hold in the full system \eqref{inc:LPR:v1:orbit}; the exponentially small
mass-deficit terms are not set to zero.
\end{proposition}
\begin{proof}
\emph{1. Finite downward variation and lower boundedness.}
Let $w_-=\max(-w,0)$. From \eqref{inc:LPR:v1:weakpair}, almost everywhere
(and in the upper Dini sense at a zero),
\[
 D^+w_-\le-\tfrac12w_-+h.
\]
Thus
\[
 w_-(s)\le e^{-s/2}w_-(0)+\int_0^s e^{-(s-t)/2}h(t)\dd t,
 \quad
 \int_0^\infty w_-\dd s\le2w_-(0)+2\int_0^\infty h\dd s<\infty.
\]
Since $f\ge0$ and $0<Q<1$, $(u')_-\le w_-/2$. Consequently $u$ has
finite total downward variation and a finite lower bound $u_-$.
The same convolution bound shows that $w_-$ is bounded. The function
$g(u)=K(u)-u$ is decreasing. Whenever $w\ge0$,
\[
 w'\le\tfrac12g(u_-)-\tfrac12w.
\]
A first-crossing comparison bounds $w_+$ by
$D_+=\max\{w_+(0),g(u_-)\}$. Combining the bounds gives $|w|\le D$.

\emph{2. The weak input tends to positive infinity.}
Suppose $u\le U$ for all $s\ge0$. Then $g(u)\ge g(U)>0$ and $h\to0$.
Choose $s_1$ such that $h\le g(U)/4$ for $s\ge s_1$, and write
$\gamma=g(U)/4$. The gap equation gives
$w'\ge\gamma-a(s)w$. For $w\le\gamma/2$ this derivative is at least
$\gamma/2$ (and at least $\gamma$ if $w\le0$). Thus $w$ reaches and
thereafter stays above $\gamma/2$ in finite time. It follows that
\[
 u'\ge\gamma\Phi(u_-)/4>0
\]
eventually, contradicting $u\le U$. Hence $u$ is unbounded above. Local
boundedness on finite intervals and the finite downward variation imply
$u(s)\to+\infty$ by \cite{LP}, \texttt{downward}. Since $|w|\le D$, also
$c=u+w\to+\infty$.

\emph{3. A tail-corrected sum gives logarithmic growth.}
Adding the first two equations of \eqref{inc:LPR:v1:weakpair} gives
\[
 (u+c)'=\tfrac12[g(u)-(1-Q)w]
       -\tfrac{r_n}{2}[\lambda u+\rho K(\rho x)].
\]
For all sufficiently late times $u\ge0$, and hence
$(u+c)'\le g(u)/2+w_-/2$.
Define the finite nonnegative tail
\[
 B(s)=\frac12\int_s^\infty w_-(t)\dd t,\qquad
 \widetilde W=u+c+B(s).
\]
Then $B\to0$ and $\widetilde W'\le g(u)/2$ at all sufficiently late
times. Since $\widetilde W=2u+w+B$,
\[
 u\ge\frac{\widetilde W-D_+-B(0)}2.
\]
Set $z=\widetilde W-D_+-B(0)$ and choose a late $s_0$ such that
$z\ge1$ for every $s\ge s_0$, possible since $u,c\to\infty$.
For nonnegative arguments $g(t)\le\phi(t)$, so
\[
 z'\le\tfrac12g(u)\le\tfrac{\phi(0)}2e^{-z^2/8}.
\]
With $F(z)=\int_0^z e^{r^2/8}\dd r$, the chain rule gives
\[
 F(z(s))\le F(z(s_0))+\tfrac{\phi(0)}2(s-s_0)=:\mathcal Q(s).
\]
For $z\ge1$, $F(z)\ge e^{(z-1)^2/8}$; therefore
$z\le1+\sqrt{8\log\mathcal Q(s)}$.
Using the lower bound $w\ge-D$ and $B\ge0$, we conclude
\begin{equation}\label{inc:LPR:v1:weakupper}
 u(s)\le A+\sqrt{2\log\mathcal Q(s)}
\end{equation}
for a finite $A>0$. Bounded $w$ gives the corresponding bound for $c$.

\emph{4. Positive gap and positive weak-input velocity.}
For $u\ge0$, restricting the Gaussian integral to $[u+1,u+2]$ gives
\[
 g(u)=\int_u^\infty(t-u)\phi(t)\dd t\ge\phi(u+2).
\]
Equation \eqref{inc:LPR:v1:weakupper} and
$(a+b)^2/2\le a^2+b^2$ show
\[
 g(u(s))\ge\phi(0)e^{-(A+2)^2}\mathcal Q(s)^{-2}
 =:\gamma_w\mathcal Q(s)^{-2}.
\]
Here $\gamma_w>0$ is a constant of this weak-coordinate estimate.
Since $h$ decays exponentially, there is $s_2\ge s_0$ such that
$h(s)\le g(u(s))/4$ for $s\ge s_2$. Thus
$w'+a(s)w\ge g(u)/4$.
Variation of constants, using $1/2\le a\le1$ and restricting its positive
integral to the last unit time interval, yields for $s\ge s_2+1$
\[
 w(s)\ge-|w(s_2)|e^{-(s-s_2)/2}
       +\frac{\gamma_w(1-e^{-1})}{4\mathcal Q(s)^2}.
\]
This is positive for all sufficiently large $s$. Then
$u'=w\Phi(u)/2+f>0$.
\end{proof}

\begin{corollary}[All four angular coordinates grow at most logarithmically]
\label{inc:LPR:v1:allgrowth}
For every fixed $0<\rho<1$ on the selected coherent orbit,
\[
 \max\{|x(s)|,|y(s)|,|u(s)|,|v(s)|\}
 =O(\sqrt{\log(2+s)}).
\]
\end{corollary}
\begin{proof}
The weak coordinates are covered by
Proposition~\ref{inc:LPR:v1:weakgrowththeorem}. The cited strong-coordinate
result gives $x=O(\sqrt{\log(2+s)})$ and $0<v-x\le D_s$; it therefore gives
the same upper-growth bound for $v$.
\end{proof}

\section{What this does and does not imply for growing horizons}
\label{inc:LPR:v1:horizons}
A coordinate chart alone is not a dynamical approximation theorem.
Corollary~\ref{inc:LPR:v1:allgrowth} implies geometrically that
\[
 q\max(|x|,|y|,|u|,|v|)\to0\quad\text{up to }s=T_q
\]
whenever $q\sqrt{\log(2+T_q)}\to0$ along this orbit (with the fixed
parameter constants understood). For example, $T_q=\exp(q^{-\beta})$
with $0<\beta<2$ meets this necessary chart-smallness condition.
It does not establish that the original flow tracks the limiting flow to
that horizon, that uniform remainder constants stay bounded, or that
background output is negligible.

If a neutral quantity accumulates an error bounded only by $Cq^2T$,
then $T\asymp q^{-2}$ is an order-one-error scale, not a vanishing-error
conclusion. For an $o(1)$ scaled error from that estimate alone, one needs
$q^2T_q\to0$, with extra logarithmic factors if $C$ grows with the chart
radius. Such drift is not automatic for a stable saturated family mass:

\begin{lemma}[A saturated mass defect can be damped]
\label{inc:LPR:v1:massdamping}
Suppose an absolutely continuous mass satisfies
\[
 N'=\lambda N(1-N)+q^2r(s),\qquad N(s)\ge n_*>0,
 \qquad |r(s)|\le C_r
\]
on $[s_0,T]$. Then
\[
 |N(s)-1|\le e^{-\lambda n_*(s-s_0)}|N(s_0)-1|
       +\frac{q^2C_r}{\lambda n_*}
          (1-e^{-\lambda n_*(s-s_0)}).
\]
\end{lemma}
\begin{proof}
Set $\delta=N-1$. Exactly $\delta'=-\lambda N\delta+q^2r$.
Thus $D^+|\delta|\le-\lambda n_*|\delta|+q^2C_r$. Integrate.
\end{proof}

\begin{remark}
The lemma assumes, and does not prove, the lower mass guard and a useful
finite-noise bound on $r$. Relative label weights can still accumulate
$q^2T$ errors when their log growth rates differ, since they do not share
the same scalar restoring mechanism. Consequently neither a certified
$T_q\asymp q^{-2}$ horizon nor a prohibition on longer horizons follows
from the current equations. A fixed-$q$ comparison to an experimental
physical horizon remains separate.
\end{remark}

\section{An exact retention identity for the original population}
\label{inc:LPR:v1:retention}
Return to the original normalized Gaussian population. Let $\lambda_0$
be the uniform initial-label law, and let
\[
 R_\tau(\theta)=\E_P[(X-f_\tau(X))X^\top
                 \mathbf1_{\{a_\theta^\top X>0\}}].
\]
The balanced characteristic mass obeys
\[
 (\log m_\alpha)'=r_\alpha(\tau),\qquad
 r_\alpha(\tau)=2b_{\psi_\alpha}^\top
                    R_\tau(\theta_\alpha)a_{\theta_\alpha}.
\]
All neurons are included in $R_\tau$. Let $C$ be a fixed measurable label
set, $N_C(\tau)=\int_Cm_\alpha(\tau)\dd\lambda_0(\alpha)$, and assume
$N_C(\tau_a)>0$.

\begin{proposition}[Fixed-label retention from a negative-growth budget]
\label{inc:LPR:v1:retentiontheorem}
For finite $\tau_a<\tau_b$, define
\[
 B_\alpha=\int_{\tau_a}^{\tau_b}[-r_\alpha(\tau)]_+\dd\tau,
 \qquad
 \overline B_C=\frac1{N_C(\tau_a)}
      \int_Cm_\alpha(\tau_a)B_\alpha\dd\lambda_0(\alpha).
\]
If this average is finite, then
\[
 \boxed{N_C(\tau_b)\ge N_C(\tau_a)e^{-\overline B_C}.}
\]
In particular a uniform labelwise bound $B_\alpha\le B$ implies
$N_C(\tau_b)\ge e^{-B}N_C(\tau_a)$.
\end{proposition}
\begin{proof}
Integrating the exact log-mass equation gives
\[
 m_\alpha(\tau_b)=m_\alpha(\tau_a)
       \exp\!\left(\int_{\tau_a}^{\tau_b}r_\alpha\dd\tau\right)
 \ge m_\alpha(\tau_a)e^{-B_\alpha}.
\]
Integrate over the fixed label set. The measure
$m_\alpha(\tau_a)\dd\lambda_0/N_C(\tau_a)$ is a probability measure.
Jensen's inequality for $e^{-z}$ gives the stated average-budget bound.
\end{proof}

\begin{remark}[Retention is not yet capture]
This elementary exact implication isolates a quantitative target. To use
it in AN2 one must prove that a fixed set of labels carries a specified
fraction of the available seed, reaches the weak chart, and has an
acceptable negative-growth budget under the full evolving residual.
No such trapping or budget is supplied by this proposition. A large mass
inside an instantaneous angular mask does not by itself prove a retained
seed. If an analytically identified captured subset has only a fraction
$\vartheta$ of the earlier mass, the lower bound is multiplied by
$\vartheta$. Mass loss through changing a mask must not be conflated with
radial loss of fixed labels.
\end{remark}

\section{Conditional exponent cancellation and the next open lemma}
\label{inc:LPR:v1:exponents}
Let $a_\rho$ solve $a_\rho=\rho K(\rho a_\rho)$ and put
\[
 P_\rho=\Phi(\rho a_\rho),\qquad
 \alpha=\frac\lambda2P_\rho.
\]
The following is a conditional algebraic implication, not an assertion
that its entry or nonlinear passage hypotheses have been established.

\begin{proposition}[Retained-seed exponent and shape compatibility]
\label{inc:LPR:v1:exponenttheorem}
Fix $0<\rho<1$, $0<S_0<1$. Suppose quantities at one common entry time
satisfy
\[
 N_{\rm ent}\ge c S_0^{1-\lambda},\qquad
 \delta_{\rm ent}\le C_s S_0^{1/2-\alpha},
\]
and a later shape measure obeys
\[
 \delta_{\rm half}\le C_p\delta_{\rm ent}
               N_{\rm ent}^{-\alpha/\lambda}+\epsilon_{\rm pass},
 \qquad \epsilon_{\rm pass}\ge0.
\]
Then
\[
 \delta_{\rm half}\le C_pC_s c^{-\alpha/\lambda}
                 S_0^{\eta(\rho)}+\epsilon_{\rm pass},\qquad
 \boxed{\eta(\rho)=\frac{1-P_\rho}{2}>0.}
\]
For a compact $\rho$ interval in $(0,1)$, this exponent has a positive
minimum. Using AN06's resident bound on $[13/20,7/10]$ gives
$\eta>97/500$ there.
More generally, seed exponent $1-\lambda+\epsilon_s+b$ and shape exponent
$1/2-\alpha-\epsilon_\delta$ produce
\[
 \eta_{\rm eff}=\frac{1-P_\rho}{2}
               -\frac{P_\rho}{2}(\epsilon_s+b)-\epsilon_\delta.
\]
\end{proposition}
\begin{proof}
Insert the two entry inequalities into the passage bound. The power is
\[
 \frac12-\alpha-(1-\lambda)\frac\alpha\lambda
 =\frac12-\frac\alpha\lambda=\frac{1-P_\rho}{2}.
\]
For fixed $\rho<1$, $a_\rho$ is finite, so $P_\rho<1$. Smooth dependence
of the resident root gives a positive minimum on compact parameter sets.
On the AN06 interval, $a_\rho<2/5$ and $\phi(0)<2/5$ give
$P_\rho<1/2+(2/5)(7/10)(2/5)=153/250$, hence
$(1-P_\rho)/2>97/500$. Including the extra exponents gives the last
identity.
\end{proof}

The exponent $b$ has a direct retention interpretation. If a captured
fraction starts with seed mass $cS_0^{1-\lambda+\epsilon_s}$ and
$\overline B_C\le B_0+b\log(1/S_0)$, then the preceding retention
proposition gives a final seed bounded below by
$c e^{-B_0}S_0^{1-\lambda+\epsilon_s+b}$, with the capture fraction
included in $c$. A bound independent of $S_0$ loses no initialization
power. A bound $o(\log(1/S_0))$ preserves the leading exponent, up to
arbitrarily small losses. Constants and errors depending on $q$ must
remain visible in any joint scaling.

\paragraph{The measurement does not replace the entry hypothesis.}
The user's reported exponent near $0.576$ was measured directly for an
instantaneous weak-family mass at strong half-mass; it was not inferred
from a response-clock regression. That is relevant evidence for the seed
scale, but its family mask and entry time must be reconciled with fixed
captured labels. AN02 currently proves a different incoming reservoir
with a conservative $S_0^{17/16}$ lower power, not a captured
$S_0^{1-\lambda}$ seed. Nor does local strong-shape linearization establish
the original-flow $S_0^{1/2-\alpha}$ upper bound. The two displayed entry
estimates must apply at the same time, without double-counting the
preceding strong-learning contraction. This file does not use an
unavailable statement denoted $G_\infty$ as an implicit assumption.

\paragraph{Next open target: a coupled retention-and-capture lemma.}
Starting from the specified isotropic initialization, prove a strong-stage
time and an analytically specified label reservoir with a quantitatively
useful mass. Then exhibit a captured fixed subset, show its entry into a
bounded weak chart under the evolving strong residual, and control its
negative integrated radial growth without an excessive initialization
power loss. The same interval must supply strong mean/shape and background
bounds at a common resident entry time. The subsequent nonlinear passage
must show that its accumulated coefficient errors fit inside the strict
exponent margin, rather than invoking positivity of the frozen exponent
alone. No numerical label mask, empirical seed law, or observed capture
may be inserted as a proved premise of the initialized theorem.

\section*{Scope and review protocol}
The combined theorem and weak-coordinate theorem are conclusions in the
selected limiting model. The retention identity is exact for the original
flow but leaves its sign-budget and capture inputs open. The exponent
calculation is conditional on unproved original-flow entry and passage
estimates. No initialized version of A or B is claimed. No canonical
non-small-bridge theorem or growing-horizon error estimate is claimed.
Deliver this file and its companion review note for independent checking;
merge nothing into the checkpoint without explicit user approval.

\begin{thebibliography}{9}
\bibitem{F} \emph{ReLU Population Foundations}, version 2.
\path{RELU_POPULATION_FOUNDATIONS_v2.tex}. Exact characteristic flow,
balance, existence, and mass identities.
\bibitem{P} \emph{Analytical Progress toward Initialized Population
Reversal}, checkpoint 1.0. \path{THEOREM_A_ANALYTICAL_PROGRESS.tex}.
\texttt{pa:branch}, \texttt{pa:limitreversal}.
\bibitem{W} \emph{Learned Limiting-Orbit Witnesses}, version 2.
\path{AN06_learned_orbit_witnesses_v2.tex}. Resident bounds and common
learned-witness rectangle.
\bibitem{LP} \emph{Linear Contrast and Limit Persistence}, version 1.
\path{LINEAR_CONTRAST_LIMIT_PERSISTENCE_v1.tex}. Eventual inactivity,
strong-coordinate growth, and compact-parameter clearing.
\bibitem{SP} \emph{Scaled Population Splitting Theory}.
\path{SCALED_POPULATION_SPLITTING_THEORY.md}. Resident shape spectra,
linear passage multipliers, and finite-noise scope distinctions.
\bibitem{A2} \emph{Initialized Weak-Side Arrival}, version 1.
\path{AN02_weak_side_arrival_v1.tex}. Incoming reservoir theorem;
capture and later retention remain open.
\end{thebibliography}
\end{document}
